\documentclass[11pt,oneside]{article}

\pdfminorversion=7
\usepackage{QuizStyle}

\hypersetup{
    pdftitle={20260914 Quiz for MATH210 Applied Complex Variables},
    pdfauthor={Jungwoo Kim},
    pdfsubject={Quiz for MATH210 Applied Complex Variables},
    pdfcreator={LaTeX}
}

\begin{document}

\quiztitle{20260914 Quiz}

\begin{exercise}{3.1}
Reduce each of these quantities to a real number:
\begin{exercisesubparts}
    \item \(\displaystyle \frac{1+2i}{3-4i}+\frac{2-i}{5i}\).
    \item \(\displaystyle \frac{5i}{(1-i)(2-i)(3-i)}\).
    \item \((1-i)^4\).
\end{exercisesubparts}
\end{exercise}

\begin{solution}
\begin{exercisesubparts}
    \item Multiplying the first denominator by its conjugate and simplifying the second term,
    \[
    \begin{aligned}
        \frac{1+2i}{3-4i}
        &=\frac{(1+2i)(3+4i)}{25}
        =-\frac15+\frac25i,\\
        \frac{2-i}{5i}
        &=\frac{(2-i)(-i)}5
        =-\frac15-\frac25i.
    \end{aligned}
    \]
    Hence
    \[
        \boxed{\frac{1+2i}{3-4i}+\frac{2-i}{5i}=-\frac25}.
    \]

    \item Since
    \[
        (1-i)(2-i)=1-3i
    \]
    and
    \[
        (1-3i)(3-i)=-10i,
    \]
    we obtain
    \[
        \boxed{\frac{5i}{(1-i)(2-i)(3-i)}=-\frac12}.
    \]

    \item We have
    \[
        (1-i)^2=-2i,
    \]
    so
    \[
        \boxed{(1-i)^4=(-2i)^2=-4}.
    \]
\end{exercisesubparts}
\end{solution}

\Needspace{12\baselineskip}
\begin{exercise}{5.4}
Verify that
\[
    \sqrt{2}\,|z|\ge |\RePart z|+|\ImPart z|.
\]
\begin{suggestion}
Reduce this inequality to \((|x|-|y|)^2\ge0\).
\end{suggestion}
\end{exercise}

\begin{solution}
Write \(z=x+iy\). Both sides are nonnegative, so it is enough to compare their squares. The desired inequality is equivalent to
\[
    2(x^2+y^2)\ge (|x|+|y|)^2,
\]
which reduces to
\[
    x^2+y^2-2|xy|=(|x|-|y|)^2\ge0.
\]
Therefore
\[
    \boxed{\sqrt{2}\,|z|\ge |\RePart z|+|\ImPart z|}.
\]
\end{solution}

\begin{exercise}{5.7}
Show that for \(R\) sufficiently large, the polynomial \(P(z)\) in Example 3, Sec.~5, satisfies the inequality
\[
    |P(z)|<2|a_n||z|^n
    \qquad\text{whenever }|z|>R.
\]
\begin{suggestion}
Observe that there is a positive number \(R\) such that the modulus of each quotient in inequality (9), Sec.~5, is less than \(|a_n|/n\) when \(|z|>R\).
\end{suggestion}
\end{exercise}

\begin{solution}
As in Example 3, Sec.~5, write
\[
    P(z)=(a_n+w)z^n,
\]
where
\[
    w=\frac{a_0}{z^n}+\frac{a_1}{z^{n-1}}+\cdots+\frac{a_{n-1}}{z}.
\]
For sufficiently large \(R\), each of the \(n\) terms on the right has modulus less than \(|a_n|/n\) whenever \(|z|>R\). Hence
\[
    |w|<|a_n|.
\]
Using the triangle inequality,
\[
    |P(z)|=|a_n+w||z|^n
    \le (|a_n|+|w|)|z|^n
    <2|a_n||z|^n.
\]
Thus the stated inequality holds for all \(|z|>R\).
\end{solution}

\begin{exercise}{6.7}
Show that
\[
    \left|\RePart\bigl(2+\overline z+z^3\bigr)\right|\le4
    \qquad\text{when }|z|\le1.
\]
\end{exercise}

\begin{solution}
For every complex number \(w\), \(|\RePart w|\le |w|\). Hence, using the triangle inequality and \(|\overline z|=|z|\),
\[
\begin{aligned}
    \left|\RePart\bigl(2+\overline z+z^3\bigr)\right|
    &\le \left|2+\overline z+z^3\right|\\
    &\le 2+|\overline z|+|z|^3\\
    &=2+|z|+|z|^3\\
    &\le4,
\end{aligned}
\]
since \(|z|\le1\).
\end{solution}

\begin{exercise}{6.13}
Show that the equation \(|z-z_0|=R\) of a circle, centered at \(z_0\) with radius \(R\), can be written
\[
    |z|^2-2\RePart\bigl(z\overline{z_0}\bigr)+|z_0|^2=R^2.
\]
\end{exercise}

\begin{solution}
Squaring \(|z-z_0|=R\) gives
\[
\begin{aligned}
    R^2
    &=|z-z_0|^2\\
    &=(z-z_0)(\overline z-\overline{z_0})\\
    &=|z|^2-z\overline{z_0}-\overline z z_0+|z_0|^2.
\end{aligned}
\]
The middle two terms are conjugates, so
\[
    z\overline{z_0}+\overline z z_0
    =2\RePart\bigl(z\overline{z_0}\bigr).
\]
Therefore
\[
    \boxed{|z|^2-2\RePart\bigl(z\overline{z_0}\bigr)+|z_0|^2=R^2}.
\]
\end{solution}

\begin{exercise}{9.9}
Establish the identity
\[
    1+z+z^2+\cdots+z^n=\frac{1-z^{n+1}}{1-z}
    \qquad(z\ne1)
\]
and then use it to derive Lagrange's trigonometric identity:
\[
    1+\cos\theta+\cos2\theta+\cdots+\cos n\theta
    =\frac12+\frac{\sin((2n+1)\theta/2)}{2\sin(\theta/2)}
    \qquad(0<\theta<2\pi).
\]
\begin{suggestion}
As for the first identity, write \(S=1+z+z^2+\cdots+z^n\) and consider the difference \(S-zS\). To derive the second identity, write \(z=e^{i\theta}\) in the first one.
\end{suggestion}
\end{exercise}

\begin{solution}
Let
\[
    S=1+z+z^2+\cdots+z^n.
\]
Then
\[
    S-zS=1-z^{n+1},
\]
so, when \(z\ne1\),
\[
    \boxed{S=\frac{1-z^{n+1}}{1-z}}.
\]
Now let \(z=e^{i\theta}\), where \(0<\theta<2\pi\). Then
\[
\begin{aligned}
    1+e^{i\theta}+\cdots+e^{in\theta}
    &=\frac{1-e^{i(n+1)\theta}}{1-e^{i\theta}}\\
    &=e^{in\theta/2}
      \frac{\sin((n+1)\theta/2)}{\sin(\theta/2)}.
\end{aligned}
\]
Taking real parts gives
\[
    1+\cos\theta+\cdots+\cos n\theta
    =\frac{\sin((n+1)\theta/2)\cos(n\theta/2)}{\sin(\theta/2)}.
\]
Using
\[
    2\sin\frac{(n+1)\theta}{2}\cos\frac{n\theta}{2}
    =\sin\frac{(2n+1)\theta}{2}+\sin\frac\theta2,
\]
we obtain
\[
    \boxed{
    1+\cos\theta+\cdots+\cos n\theta
    =\frac12+\frac{\sin((2n+1)\theta/2)}{2\sin(\theta/2)}
    }.
\]
\end{solution}

\begin{exercise}{11.6}
Find the four zeros of the polynomial \(z^4+4\), one of them being
\[
    z_0=\sqrt2\,e^{i\pi/4}=1+i.
\]
Then use those zeros to factor \(z^2+4\) into quadratic factors with real coefficients.
\end{exercise}

\begin{solution}
The equation \(z^4+4=0\) is
\[
    z^4=4e^{i(\pi+2k\pi)}.
\]
Thus its four roots are
\[
    z=\sqrt2\,e^{i(\pi/4+k\pi/2)},
    \qquad k=0,1,2,3,
\]
namely
\[
    1+i,\quad -1+i,\quad -1-i,\quad 1-i.
\]
Pairing conjugate roots gives real quadratic factors:
\[
\begin{aligned}
    z^4+4
    &=(z-(1+i))(z-(1-i))(z-(-1+i))(z-(-1-i))\\
    &=(z^2-2z+2)(z^2+2z+2).
\end{aligned}
\]
Hence
\[
    \boxed{z^4+4=(z^2+2z+2)(z^2-2z+2)}.
\]
\end{solution}

\begin{remark}
The second polynomial in the printed statement is a typographical error. The four zeros requested are the zeros of \(z^4+4\), and the printed answer factors \(z^4+4\). Accordingly, the intended factorization is that of \(z^4+4\).
\end{remark}

\begin{exercise}{12.7}
Determine the accumulation points of each of the following sets:
\begin{exercisesubparts}
    \item \(z_n=i^n\), \(n=1,2,\ldots\).
    \item \(z_n=i^n/n\), \(n=1,2,\ldots\).
    \item \(0\le \arg z<\pi/2\), \(z\ne0\).
    \item \(\displaystyle z_n=(-1)^n(1+i)\frac{n-1}{n}\), \(n=1,2,\ldots\).
\end{exercisesubparts}
\end{exercise}

\begin{solution}
\begin{exercisesubparts}
    \item The set consists only of the four points \(1,i,-1,-i\). A finite set has no accumulation points. Thus there are none.

    \item Since \(|z_n|=1/n\to0\), every neighborhood of \(0\) contains points of the set other than \(0\). Conversely, fix \(a\ne0\) and choose \(N\) so large that \(|z_n|<|a|/2\) for all \(n>N\). A sufficiently small deleted neighborhood of \(a\) excludes these terms and also the finitely many distinct points among \(z_1,\ldots,z_N\). Hence the only accumulation point is
    \[
        \boxed{0}.
    \]

    \item The set is the sector consisting of the nonzero points with angle from \(0\) up to, but not including, \(\pi/2\). Every point of the closed first quadrant can be approached by points of this sector, and no point outside that closed quadrant can. Thus the set of accumulation points is
    \[
        \boxed{\{x+iy\mid x\ge0,\ y\ge0\}}.
    \]

    \item For even \(n\),
    \[
        z_n=(1+i)\frac{n-1}{n}\longrightarrow 1+i,
    \]
    while for odd \(n\),
    \[
        z_n=-(1+i)\frac{n-1}{n}\longrightarrow -(1+i).
    \]
    These are the only two accumulation points, so
    \[
        \boxed{\pm(1+i)}.
    \]
\end{exercisesubparts}
\end{solution}

\begin{exercise}{14.2}
In each case, write the function \(f(z)\) in the form \(f(z)=u(x,y)+iv(x,y)\):
\begin{exercisesubparts}
    \item \(f(z)=z^3+z+1\).
    \item \(\displaystyle f(z)=\frac{\overline z^{\,2}}{z}\), \(z\ne0\).
\end{exercisesubparts}
\begin{suggestion}
In part (b), start by multiplying the numerator and denominator by \(\overline z\).
\end{suggestion}
\end{exercise}

\begin{solution}
Write \(z=x+iy\).
\begin{exercisesubparts}
    \item Expanding \((x+iy)^3\),
    \[
    \begin{aligned}
        f(z)
        &=(x+iy)^3+(x+iy)+1\\
        &=(x^3-3xy^2+x+1)+i(3x^2y-y^3+y).
    \end{aligned}
    \]
    Therefore
    \[
        \boxed{u=x^3-3xy^2+x+1,
        \qquad v=3x^2y-y^3+y}.
    \]

    \item Since \(z\ne0\),
    \[
        \frac{\overline z^{\,2}}{z}
        =\frac{\overline z^{\,3}}{|z|^2}.
    \]
    Now
    \[
        \overline z^{\,3}
        =(x-iy)^3
        =(x^3-3xy^2)+i(y^3-3x^2y),
    \]
    and hence
    \[
        \boxed{
        f(z)=\frac{x^3-3xy^2}{x^2+y^2}
        +i\frac{y^3-3x^2y}{x^2+y^2}
        }.
    \]
\end{exercisesubparts}
\end{solution}

\Needspace{17\baselineskip}
\begin{exercise}{14.3}
Suppose that
\[
    f(z)=x^2-y^2-2y+i(2x-2xy),
    \qquad z=x+iy.
\]
Use the expressions
\[
    x=\frac{z+\overline z}{2}
    \qquad\text{and}\qquad
    y=\frac{z-\overline z}{2i}
\]
to write \(f(z)\) in terms of \(z\), and simplify the result.
\end{exercise}

\begin{solution}
Using
\[
    x=\frac{z+\overline z}{2},
    \qquad
    y=\frac{z-\overline z}{2i},
\]
we obtain
\[
    x^2-y^2=\frac{z^2+\overline z^{\,2}}2,
    \qquad
    -2y=i(z-\overline z),
\]
and
\[
\begin{aligned}
    i(2x-2xy)
    &=i(z+\overline z)-\frac{z^2-\overline z^{\,2}}2.
\end{aligned}
\]
Therefore
\[
\begin{aligned}
    f(z)
    &=\frac{z^2+\overline z^{\,2}}2+i(z-\overline z)
      +i(z+\overline z)-\frac{z^2-\overline z^{\,2}}2\\
    &=\overline z^{\,2}+2iz.
\end{aligned}
\]
Thus, for the function as printed,
\[
    \boxed{f(z)=\overline z^{\,2}+2iz}.
\]
\end{solution}

\begin{remark}
The printed answer in the textbook is \(z^2+2iz\), which does not agree with the printed function. The direct algebra above gives \(\overline z^{\,2}+2iz\).
\end{remark}

\begin{exercise}{14.4}
Write the function
\[
    f(z)=z+\frac1z
    \qquad(z\ne0)
\]
in the form \(f(z)=u(r,\theta)+iv(r,\theta)\).
\end{exercise}

\begin{solution}
With \(z=re^{i\theta}\),
\[
    \frac1z=\frac1r e^{-i\theta}.
\]
Hence
\[
\begin{aligned}
    f(z)
    &=r(\cos\theta+i\sin\theta)
      +\frac1r(\cos\theta-i\sin\theta)\\
    &=\left(r+\frac1r\right)\cos\theta
      +i\left(r-\frac1r\right)\sin\theta.
\end{aligned}
\]
Therefore
\[
    \boxed{
    u(r,\theta)=\left(r+\frac1r\right)\cos\theta,
    \qquad
    v(r,\theta)=\left(r-\frac1r\right)\sin\theta
    }.
\]
\end{solution}

\end{document}
